\documentclass[11pt]{article}
\usepackage[margin=1in]{geometry}
\usepackage{amsmath,amsthm,amssymb}
\usepackage{booktabs,hyperref,enumitem}

\newtheorem{theorem}{Theorem}[section]
\newtheorem{lemma}[theorem]{Lemma}
\newtheorem{proposition}[theorem]{Proposition}
\newtheorem{remark}[theorem]{Remark}
\newtheorem{definition}[theorem]{Definition}

\newcommand{\eps}{\varepsilon}
\newcommand{\R}{\mathbb{R}}
\newcommand{\norm}[2]{\|#1\|_{#2}}

\title{Phi-renormalization continued:\\
corrected five-dimensional energy and the uniform-strain obstruction}
\author{Jonathan R.\ Simons\\
Prime Field Technologies LLC, Savannah, Georgia}
\date{August 2026 \quad|\quad Continuation of Zenodo 20405404}

\begin{document}
\maketitle

\begin{abstract}
This note continues the author's May 2026 Zenodo preprint
(concept DOI \href{https://doi.org/10.5281/zenodo.20405404}{10.5281/zenodo.20405404}).
That deposit introduced the \(\Phi\)-renormalization
\(\Gamma=ru_\theta\), \(\Phi=\Gamma/r^2=u_\theta/r\) and the algebraic identity
\(\frac1{r^4}\partial_z(\Gamma^2)=\partial_z(\Phi^2)\).
The present version does three things: (i)~it keeps that identity as the
author's public antecedent; (ii)~it removes reliance on an unsupported
\(\Phi\) maximum-principle / Sobolev bound; (iii)~it writes the correct
lifted equation and \(r^3\,dr\,dz\) energy for \(F=u_\theta/r\), and
isolates the remaining obstruction
\[
  \sup_{\eps\in(0,1]}\int_0^T \norm{u_r^\eps/r}{L^\infty}\,dt<\infty
\]
without assuming the desired regularity.
\textbf{No claim is made that the classical unaugmented axisymmetric-with-swirl
Navier--Stokes system is globally regular.}
Shahmurov's 2026 announcements are cited and not used.
\end{abstract}

\section{Public record}

The May 2026 files are
\href{https://doi.org/10.5281/zenodo.20405405}{10.5281/zenodo.20405405}
(\texttt{PhiRenorm\_TrackB.pdf}) and
\href{https://doi.org/10.5281/zenodo.20405597}{10.5281/zenodo.20405597}.
The algebraic identity appeared there. The later mathematics of this note
did not. A linked Zenodo ``New version'' preserves that chronology.

\section{What is proved (algebra)}

\begin{theorem}[Phi-renormalization identity]
Let \(\Gamma=ru_\theta\) and \(\Phi=u_\theta/r\) on \(\{r>0\}\). Then
\[
  \frac1{r^4}\partial_z(\Gamma^2)=\partial_z(\Phi^2).
\]
If \(\Gamma^\eps=r^2\Phi^\eps\) likewise,
\(\frac1{r^4}\partial_z(\Gamma_\eps^2-\Gamma^2)
=\partial_z\bigl((\Phi^\eps+\Phi)(\Phi^\eps-\Phi)\bigr)\).
\end{theorem}

\begin{proof}
\(\Gamma=r^2\Phi\), so \(\Gamma^2=r^4\Phi^2\). Since \(\partial_z(r^4)=0\),
\(\frac1{r^4}\partial_z(\Gamma^2)=\partial_z(\Phi^2)\).
The difference identity is the same factorization.
\end{proof}

This is a change of dependent variable. It does not, by itself, bound
\(\Phi\) in \(L^\infty\) or close a Beale--Kato--Majda integral.

\begin{remark}[Circular Sobolev line, withdrawn as a proof step]
The May remark \(\norm{\partial_z(\Phi^2)}{L^2}\le 2\norm{\Phi}{\infty}\norm{\partial_z\Phi}{L^2}\)
``by Sobolev on a bounded domain'' assumes the \(L^\infty\) control that
classical continuation is trying to prove. It is withdrawn as a load-bearing step
for the unaugmented system.
\end{remark}

\section{Corrected \(F\)-equation and five-dimensional energy}

Write \(F:=u_\theta/r\) (the same intensive field as \(\Phi\)).
The standard axisymmetric calculus gives, for smooth solutions,
\begin{equation}\label{eq:F}
\partial_t F+u_r\partial_r F+u_z\partial_z F+2\frac{u_r}{r}F
=\nu\Bigl(\partial_{rr}+\frac3r\partial_r+\partial_{zz}\Bigr)F.
\end{equation}
The extra \(-F/r^2\) damping belongs to the equation for \(u_\theta\), not for \(F\).
The natural energy measure for~\eqref{eq:F} is the five-dimensional lift
\[
  d\mu_5:= r^3\,dr\,dz,
\]
not the three-dimensional weight \(r\,dr\,dz\).

\begin{proposition}[Energy identity with strain pairing]
For smooth compactly supported axisymmetric fields, testing~\eqref{eq:F}
against \(F\,d\mu_5\) yields
\begin{equation}\label{eq:energy}
  \frac12\frac{d}{dt}\int F^2\,d\mu_5
  +\nu\int |\nabla_5 F|^2\,d\mu_5
  =-2\int \frac{u_r}{r}\,F^2\,d\mu_5,
\end{equation}
up to the usual integration-by-parts justification of the five-dimensional
viscous form. The right-hand side is the strain pairing.
\end{proposition}

\begin{proof}[Proof outline]
The transport field \((u_r,u_z)\) is divergence-free in the axisymmetric
meridional plane after the \(r\)-weight of incompressibility is used;
the remaining first-order contribution is the stretching
\(2(u_r/r)F\). Multiplying by \(F\) and integrating in \(d\mu_5\) produces
the displayed pairing. The viscous operator on the right of~\eqref{eq:F}
is (minus) the Dirichlet form of the five-dimensional Laplacian in the lift.
\end{proof}

\section{The obstruction}

\begin{proposition}[Continuation if the strain is integrable]
If a smooth axisymmetric solution on \([0,T_*)\) satisfies
\begin{equation}\label{eq:target}
  \int_0^{T} \norm{u_r/r}{L^\infty}\,dt<\infty
  \qquad\text{for every }T<T_*,
\end{equation}
then \(\int F^2\,d\mu_5\) remains finite on \([0,T]\) by Gronwall
on~\eqref{eq:energy}, and the standard swirl continuation criteria apply.
\end{proposition}

The target for a classical global-regularity proof is therefore to obtain
\eqref{eq:target} \emph{without assuming}~\eqref{eq:target}.
A family of regularizations \(u^\eps\) may be used only if the bound is
uniform in \(\eps\in(0,1]\).

\begin{remark}[Why a cubic energy does not close]
An inequality of the form \(E'+c\nu D\le C\nu^{-3}E^3\) is supercritical.
The comparison ODE \(y'=C\nu^{-3}y^3\) explodes in finite time for large \(y(0)\).
Such an estimate is consistent with both global existence and blowup; it is
not a global bound.
\end{remark}

\begin{remark}[Fractional bookkeeping]
Testing a fractional Laplacian of order \(\sigma\) produces an energy in
\(\dot H^{\sigma}\), not \(\dot H^{2\sigma}\).
\end{remark}

\section{What is not claimed}

\begin{itemize}[leftmargin=1.2em]
\item Classical (unaugmented) large-data global regularity in the
axisymmetric swirl class.
\item A maximum principle for \(\Phi\) or \(F\) beyond the classical
maximum principle for \(\Gamma=ru_\theta\).
\item Gronwall-free convergence of a \(Q_1\)-approximation to a
classical solution on a maximal interval that may be finite.
\item Any prime-lattice, \(Q_6\), or spectral-clock equivalence.
\end{itemize}

A \(Q_1\)-augmented system with uniformly positive extra viscosity is a
different PDE. Global smoothness for that PDE, if written carefully, is a
method note. It is not the classical swirl theorem.

\section{Related announcements}

Shahmurov has released several 2026 preprints on axisymmetric swirl and
on full-system Navier--Stokes, including
arXiv:2604.03519, arXiv:2605.01875, and arXiv:2606.07869.
Those papers use, among other devices, an \((A,W)\) circulation-gradient
pair and an axis Hardy formula for \(\Gamma\).
The present note does not employ that construction.
If a later independent proof of~\eqref{eq:target} is obtained on the
variables of~\S3, it may be stated as independent, with the May 2026
Zenodo record as antecedent for the \(\Phi\)-identity only.

\section{Next checkable steps}

\begin{enumerate}[leftmargin=1.2em]
\item Write~\eqref{eq:energy} with every boundary term and cutoff error
displayed (no ``up to the usual'').
\item Attempt a subcritical bound on \(\int (u_r/r)F^2\,d\mu_5\) from
\(\Gamma\in L^\infty\) (classical) and the five-dimensional dissipation
alone --- or exhibit the first scaling obstruction.
\item Independently audit the first Hardy/endpoint lemma in
arXiv:2606.07869. If it fails, record the exact line. If it holds,
cite it and keep this route separate.
\end{enumerate}

\begin{thebibliography}{99}

\bibitem{SimonsPhi2026}
J.\ R.\ Simons,
\textit{Phi-Renormalization and Global Regularity for Axisymmetric-with-Swirl Navier--Stokes},
Zenodo, 27 May 2026.
\href{https://doi.org/10.5281/zenodo.20405404}{10.5281/zenodo.20405404}.

\bibitem{ShahApr}
R.\ Shahmurov,
\textit{Unconditional Axis-Regularity in the 5D Corridor},
arXiv:2604.03519.

\bibitem{ShahMay}
R.\ Shahmurov,
\textit{Large-Data Global Regularity for Three-Dimensional Navier--Stokes I},
arXiv:2605.01875.

\bibitem{ShahJun}
R.\ Shahmurov,
\textit{Global Regularity for Axisymmetric Navier--Stokes Flows with Swirl},
arXiv:2606.07869.

\bibitem{Lady1968}
O.\ A.\ Ladyzhenskaya,
\textit{Unique global solvability of the three-dimensional Cauchy problem
for the NS equations in the presence of axial symmetry},
Zap.\ Nauchn.\ Sem.\ LOMI \textbf{7} (1968), 155--177.

\end{thebibliography}

\end{document}
